\chapter{Problem Formulation}

This thesis describes an autonomous robot executing an active object search task in an intelligent way. Active object search is the problem of finding an efficient policy to find an object in an large scale, unknown, $3$D indoor environment $\psi$, with $\psi \subseteq \mathbb{R}^3$. 

In this thesis an object is defined as $O := (x,y,z,o)$, where the vector $(x,y,z)$ is the position of the object and $o$ is the object class or object type with $o \in C_{O}$ and $C_{O}$ being the set of detectable object classes. 

An indoor environment can be segmented into a set of rooms $R := \{r_{1}, r_{2}, ..., r_{K}\}$, with K being the number of rooms and a room $r_i \subseteq \mathbb{R}^3$. Further holds $\psi = r_{1} \cup r_{2} \cup ... \cup r_{K}$ and $r_{k} \cap r_{l} = \emptyset$ $\forall k \in K, l \in K, k \neq l$. 

A set of probability distributions $P_{\psi}(x,y,z) = \{P_{\psi}^{o_{1}}(x,y,z),..., P_{\psi}^{o_{\mid C_O\mid}}(x,y,z)\}$ can be introduced describing the probability of an object of type $o_i$ being located at position $(x,y,z)$ $P_{\psi}^{o_{i}}(x,y,z)$. In contrast to an uninformed object search these probability distributions are not flat enabling the possibility to find a more intelligent search policy.

The probability of an object being located at a certain position is correlated with the appearance of the surrounding space. In this thesis the appearance is classified into a room type $rt$ with $rt \in C_{rt}$ and $C_{rt}$ being the set of different room types. Within a room the appearance is similar or there a small number of regions in the room with similar appearance. 

A search task is searching for an object of a given target object class $o_i \in C_{O}$ in the current indoor environment given the prior knowledge of $\psi$ and $P_{\psi}^{o_{i}}(x,y,z)$ until it is found or another termination criteria (e.g. the whole environment was searched) is met. This is different to other research in this area, where object instances were searched. In this thesis the prior knowledge about $\psi$ and $P_{\psi}^{o_{i}}(x,y,z)$ is zero or incomplete.

The robot can execute two dimensional motion actions to goal pose $g := (x,y,\theta)$ in the motion space $\phi$, with $\phi \in \mathbb{R}^2$. The motion space $\phi$ is a projection of the search space $\psi$ onto the xy-plane. The robot is running sensing actions periodically, where $S $ is the set of all acquired sensor information. The system should execute the search task on the path $\mathbf{p^{\ast}} := \argminA_{\mathbf{p_{i}}} (cost(\mathbf{p_i}))$ with $\mathbf{p_i}$ being the $i^{th}$ possible path $\mathbf{p_i} := (g_i^1, g_i^2, ..., g_i^{end})$ and $cost(\mathbf{p_i})$ being the cost of the path $\mathbf{p_i}$. The cost function used in this thesis is the expected search time $E(t|\mathbf{p_i}) = \int_{T_{0}}^{\infty} t \cdot f_{o_{i}}(t|\mathbf{p_i}) dt$ with $f_{o_{i}}(t)$ being the probability density function of stopping the search of object $o_{i}$ over time $t$ and $T_{0}$ being the search start time.

The probability density function $f_{o_{i}}(t|\mathbf{p_i})$ is based on the path $\mathbf{p_{i}}$ and the probability distribution $P_{\psi}^{o_{i}}(x,y,z)$. $P_{\psi}^{o_{i}}(x,y,z)$ has to be estimated based on the acquired sensor information $S$, commonsense knowledge and prior estimates. The sensor information is range data from a LIDAR and RGBD-images, the commonsense knowledge is knowledge about object appearance, room type appearance and their correlation and the prior estimates are estimates about unexplored space.

% In addition to the search space $\psi$ the probability distribution of an object of the target object class in the search space  is important for the search strategy. 
% 
% A search task is searching for an object of a given target object class $o \in C_{O}$ in the current indoor environment given the prior knowledge of $\psi$ and $P_{\psi}(o)$ until it is found or another termination criteria (e.g. the whole environment was searched) is met. This is different to other research in this area, where object instances were searched. In this thesis the prior knowledge about $\psi$ and $P_{\psi}(o)$ is zero or incomplete. \newline
% The robot can execute two dimensional motion actions to goal positions $g := (x,y,\theta)$ in the search space $\psi$ and is running sensing actions periodically, where $S $ is the set of all acquired sensor information. The system should execute the search task on the path $p^{\ast} := \argminA_{p_{i}} cost(p_i )$ with $p_i$ being the $i^{th}$ possible path $p_i := (g_i^1, g_i^2, ..., g_i^{end})$ and $cost(p_i)$ being the cost of the path $p_i$. The main cost function used in this thesis is the expected search time $E(t|p_i) = \int_{T_{0}}^{\infty} t \cdot f_{o}(t|p_i) dt$ with $f_{o}(t)$ being the probability density function of finding the object $o$ over time $t$ and $T_{0}$ being the search start time.
% 
% The probability density function $f_{o}(t|p_i)$ is based on the path $p_i$ and the probability distribution $P_{\psi}(o)$. $P_{\psi}(o)$ has to be estimated based on the acquired sensor information $S$ and additional commonsense knowledge provided to the robot, like the concept of a room, room types and correspondences of room type and object probability. The commonsense knowledge is also useful to improve necessary heuristics, which are inevitable due to the complexity of the problem. 

%The robot is equipped with commonsense knowledge about indoor environments. The most important concept the %robot is aware of is a room. The concept of rooms leads to a high-level representation of indoor environment as %a topological map with rooms as nodes and doorways as edges defined as $H := (R,D)$. \newline
%R is a set of rooms and defined with $R := \{r_{1}, r_{2}, ..., r_{K}\}$ with K being the number of rooms. %Further holds $E = r_{1} \cup r_{2} \cup ... \cup r_{K}$ and $r_{k} \cap r_{l} = \emptyset$ $\forall k \in K, l %\in K, k \neq l$. \newline
%D is a set of doorways and every doorway is connecting two rooms and is described with $d := (r_{1},r_{2})$, %where $r_{1}$ and $r_{2}$ specify the connected rooms.

%Rooms can be assigned to one or a small number of categories and the correlation between room category and %object probability is high. Therefore a reasonable approach is to split the problem into a high-level part, %operating on the topological map, and a low-level part, the search within a room.

%\section{High-level search}

%For the high-level search a room is defined as $r_k := (P(o),p_o(t),T_M)$. Where $P(o)$ is the probability of %an object of class $o$ being in the room, $p_o(t)$ is the probability distribution of ending the search at time %$t$ and $T_M := \bigcup\limits_{i=1}^{L} \bigcup\limits_{j=1}^{L} T_{M_{ij}}$ is the matrix of movement times %between 

%As the prior knowledge is incomplete the robot needs additional information. Commonsense knowledge like the concept of a room and room-type object correlations are provided to the robot to enable more intelligent and therefore more efficient search strategies.

%The robot can execute two dimensional movement actions to a goal pose $g := (x,y,\theta)$ in the search space $\phi$ and execute sensing actions at a nearly real-time rate.

%The application of commonsense knowledge enables a more intelligent and efficient search planning. One important %concept is the concept a room.  
%Indoor environment $\psi$ is a bounded three dimensional space which can be represented as a graph with rooms as %nodes and doorways as edges defined as $H := (R,D)$. 

%R is a set of rooms and defined with $R := \{r_{1}, r_{2}, ..., r_{K}\}$ with K being the number of rooms. %Further holds $E = r_{1} \cup r_{2} \cup ... \cup r_{K}$ and $r_{k} \cap r_{l} = \emptyset$ $\forall k \in K, l %\in K, k \neq l$. \newline
%D is a set of doorways and every doorway is connecting two rooms and is described with $d := %(r_{1},r_{2},P_{1},P_{2})$. Where $r_{1}$ and $r_{2}$ specify the connected rooms and $P_{1}$ and $P_{2}$ are %the poses of the doorway in the corresponding room coordinate frames with $P := (x,y,\theta)$.

%To execute a search task efficiently a more detailed representation of the indoor environment is needed. The %robot has to be able to generate accessible view poses and also has to be able to find the optimal accessible %view pose. \newline
%Therefore a room $r_{k}$ is described with $r_{k} := (map3D_{k}, mapO_{k})$. $map3D_{k}$ is a three dimensional %occupancy grid map of the room enabling navigation. $mapO_{k}$ is a set of three dimensional grid maps $mapO_{k} %:= \{mapO_{k}^1,...,mapO_{k}^{C_{O}}\}$, where $mapO_{k}^o$ specifies for every grid cell the probability of an %object of class $o$ being in this cell. 

%The main cost function used in this thesis is the expected search time defined as $E(t) = \int_{T_{0}}^{\infty} %t \cdot p(t) dt$ with $p(t)$ being the probability density function of finding the object over time and $T_{0}$ %being the search start time.
